\section{Use the chapter package for feedback and submission}
The chapter package is \texttt{033-png}. It includes the native reading fragments,
all six construction stages, input generators, public completed-interface checks,
and a separate assembly file containing only grading storage. It includes no
completed decoder or quiz answer keys. Guided work follows the one hundred native
steps. Free-writing work may use a different internal parser, table representation,
vector strategy or private frame while keeping the callable interfaces.

For guided work, begin with the loader source in
\texttt{construction/png-construction/main.asm}. Keep your working source between
observations. At each later stage, use its attachment command to add the next
unfinished regions while retaining your own earlier code. For free writing,
\texttt{starter/main.asm} contains all unfinished stages together. It is a convenient
interface scaffold, not a required implementation layout. The native chapter
remains the authority for the byte windows, permitted subset and failure behavior.

The public feedback domains are loading, collection, inflation, reconstruction,
packing, complete decoding, repeated calls and the actual caller. Run one with
\texttt{python3 grade.py working.asm --domain decode --cli /path/to/cli.jar
  --java /path/to/java}. Omit \texttt{--domain} to run all completed domains. The
runner supplies \texttt{checks/harness.asm} as another assembly source. That file
provides disjoint test buffers in the separate wide-data window; it contains no
algorithm. Do not copy it into your program or reuse its reserved
\texttt{pngCheck\_} labels. The ordinary image viewer still uses the caller's
classroom heap allocation. The checks set target byte order before assembly and
compare numeric pixel values in either order.

An early partial implementation is expected to fail completed-stage checks. Reach
the step's named checkpoint and compare its native observation while developing
that stage; run its completed domain after filling the stage. A checkpoint pass
establishes that illustrated state, not an entire parser contract. Instructor-private
checks add cases and cumulative observations. Remote delivery and automatic
progression/hint enforcement remain future work; the package records the intended
controls without claiming that the current plugin enforces them.

Submit the working assembly and the observations requested by the course. Preserve
full source files and a small record of the input, entry, selected byte order,
status/count and expected pixels or retained state. Public feedback uses known
coordinate-generated images and explicitly malformed cases, rather than the private
assembly implementation. Passing its finite cases is useful feedback, not proof
that every possible image is accepted or rejected correctly.

For a rejected image, report which layer first establishes the defect. Scratch
changes after failed collection or inflation are permitted inside their selected
windows. Framebuffer and final-context publication have a stronger contract:
completed failure must not store to them. Keep actual earlier successful output
through a rejected call, then accept a new valid image in the same execution session.
Instruction counts are measurements with named boundaries, not requirements imposed
on an equivalent correct solution.

The package includes actual IntelliJ Light figures for loading, container metadata,
reconstructed rows, vector packing and final publication. Native callouts and
retained capture receipts accompany the original screenshots. The separate
authoring plugin depends on MARS and performs the captures. Native PDF tagging and
embedded fonts are checked during composition; independent accessibility review
remains separate from simulator tests. Substantive editing is a later goal.
